% !TeX program = lualatex
% Compile twice: lualatex 183.tex
% All references and prose are embedded; no BibTeX database or external inputs.
\documentclass[11pt,a4paper]{article}
\usepackage[margin=24mm,headheight=14pt]{geometry}
\usepackage{fontspec}
\setmainfont{Latin Modern Roman}
\setsansfont{Latin Modern Sans}
\setmonofont{Latin Modern Mono}
\usepackage{amsmath,amssymb,mathtools}
\usepackage{unicode-math}
\setmathfont{Latin Modern Math}
\usepackage{microtype}
\usepackage{enumitem,xurl,needspace}
\usepackage[dvipsnames]{xcolor}
\usepackage{hyperref}
\hypersetup{unicode=true,colorlinks=true,linkcolor=MidnightBlue,urlcolor=MidnightBlue,
 pdftitle={500 Mathematical Problem Statements},pdfauthor={Source-annotated compilation},bookmarksnumbered=true}
\usepackage{bookmark}
\usepackage{tocloft}
\setlength{\cftsecnumwidth}{3em}
\cftsetpnumwidth{2.5em}
\cftsetrmarg{4em}
\usepackage{titlesec}
\titleformat{\section}{\Large\bfseries\raggedright}{\thesection}{0.7em}{}
\usepackage{fancyhdr}
\pagestyle{fancy}\fancyhf{}
\fancyhead[L]{\small\sffamily Mathematical problem statements}
\fancyhead[R]{\small Source edition 22}
\fancyfoot[C]{\thepage}
\setlength{\parindent}{0pt}
\setlength{\parskip}{5pt plus 1pt}
\setlength{\emergencystretch}{3em}
\setcounter{tocdepth}{1}
\setcounter{secnumdepth}{1}
\setlist[itemize]{leftmargin=3em,itemsep=5pt,topsep=4pt,parsep=0pt}
\allowdisplaybreaks
\newcommand{\N}{\mathbb{N}}
\newcommand{\Z}{\mathbb{Z}}
\newcommand{\Q}{\mathbb{Q}}
\newcommand{\R}{\mathbb{R}}
\newcommand{\C}{\mathbb{C}}
\newcommand{\F}{\mathbb{F}}
\newcommand{\A}{\mathbb{A}}
\newcommand{\PP}{\mathbb P}
\newcommand{\E}{\mathbb{E}}
\newcommand{\eps}{\varepsilon}
\newcommand{\dd}{\,\mathrm d}
\newcommand{\sref}[2]{\hyperlink{source:#1:#2}{[S#2]}}
\newcommand{\eref}[2]{\hyperlink{extra:#1:#2}{[E#2]}}
% Turn the next switch off for a shorter reading copy. The quotations preserve
% every exactTarget field, including qualifications omitted from a short summary.
\newif\ifshowcatalogue
\showcataloguetrue
\newcommand{\cataloguescope}[1]{\ifshowcatalogue
 \paragraph{Exact scope in the supplied catalogue.}
 \begin{quote}\small #1\end{quote}\fi}

% Standalone notebook wrapper; original problem section retained verbatim.
\numberwithin{equation}{section}
\hypersetup{pdftitle={Research attempt -- problem 183},
 pdfauthor={Research notebook; source compilation retained}}
\fancyhead[L]{\small\sffamily Research notebook}
\fancyhead[R]{\small\sffamily Problem 183 | 27 September 2026}
\begin{document}
\setcounter{section}{182}
% ================================================================
% problemId: problem.furstenberg-times-two-times-three-invariant-measures
% source releaseRank: 183; source releaseStatus: open
% exactTarget SHA-256: c8aee98394b341cdf7dac54c5e0b8a65f4f8d4f6356dfe78e39b2ba2afe3f7cc
\clearpage
\section{\texorpdfstring{Furstenberg's x2 x3 conjecture on invariant measures}{Furstenberg's x2 x3 conjecture on invariant measures}}\label{183}
\subsection{Definitions and mathematical statement}
On $\mathbb T=\R/\Z$, let $T_2(x)=2x$ and $T_3(x)=3x$. A Borel probability measure $\mu$ is invariant if $(T_2)_*\mu=(T_3)_*\mu=\mu$. Joint ergodicity means every measurable set invariant modulo null sets under both maps has measure zero or one. The selected assertion is
\[
 \mu\text{ invariant and jointly ergodic}
 \quad\Longrightarrow\quad
 \mu=\text{Haar measure or }|\operatorname{supp}\mu|<\infty.
\]
Ergodicity for the joint action must not be silently strengthened to ergodicity for each map separately. The catalogue also mentions a coprime $(a,b)$ formulation; that broader parameter statement is recorded in its quotation, but an equivalence with the single pair $(2,3)$ is not established merely by these definitions.
\par\smallskip\textit{Formulation and concept references:} \sref{183}{1}.
\cataloguescope{Prove that every ergodic probability measure on the unit circle T that is invariant under both maps x -> 2x mod 1 and x -> 3x mod 1 is either the Lebesgue measure or supported on a finite set (equivalently, for coprime a,b >= 2, every ergodic x(a,b)-invariant measure is Lebesgue or finitely supported).}
\subsection{Short English statement}
Must a jointly ergodic circle measure preserved by both doubling and tripling be either uniform or concentrated on finitely many points?
% BEGIN RESEARCH ADDITION ATTEMPT-183
\subsection{Research attempt}

\paragraph{Current frontier.}
Rudolph's positive-entropy theorem uses joint ergodicity, exactly as in
this entry: a measure invariant under a coprime pair and ergodic for its
semigroup is Lebesgue if either generator has positive entropy
\eref{183}{1}, abstract and Sections 1 and 3. Thus the atomless zero-entropy
case is the missing one, not a case excluded by the definition. The 2024
Burton--Panangaden paper supplies equivalent complex-analytic and
operator-algebraic formulations; it does not prove the classification or
its separately stated periodic-approximation conjecture
\sref{183}{1}, Section 1.1. The argument below tests a Fourier-energy
approach directly on the measure, without strengthening joint ergodicity
to ergodicity for each map.

\paragraph{Attack route.}
An entropy route needs to force positive entropy in a nonatomic case where
zero entropy is currently permitted. A periodic-approximation route needs
both approximation and control of the limiting periodic measures. The
selected route uses positive definiteness and the exact replication of
Fourier coefficients under multiplication by two and three. The hoped-for
input is an upper bound on low-frequency Fourier energy sharp enough to
contradict that replication. We derive the required scale, prove why
nonatomicity alone is far too weak, and test a simple coefficient-amplifying
inequality.

\paragraph{Attempt: completely dispose of the atomic case.}
Suppose an invariant jointly ergodic probability measure $\mu$ has an atom.
Its largest atom mass $p>0$ is attained: only finitely many atoms can have
mass at least half the positive supremum. Let $S$ be the nonempty finite
set of atoms of mass $p$. For $T=T_2$ or $T_3$, invariance implies
$\mu(\{Tx\})\ge\mu(\{x\})=p$ when $x\in S$, so $T(S)\subseteq S$.
Two points of $S$ cannot have the same image, since its mass would then be
at least $2p$. Thus each $T$ permutes $S$. Moreover $S\subseteq T^{-1}S$
and $\mu(T^{-1}S)=\mu(S)$, so $S$ is invariant modulo null sets.
Joint ergodicity gives $\mu(S)=1$.

The measure is consequently uniform on a single joint orbit in $S$. Each
point is periodic under both maps separately, so its reduced rational
denominator divides some $2^r-1$ and some $3^s-1$, and is coprime to six.
These conclusions were proved using joint ergodicity only. They agree
with the finite-support analysis in \sref{183}{1}, Section 2.6. The
unresolved part is therefore exactly the nonatomic case.

\paragraph{Attempt: a logarithmic-square Fourier obstruction.}
Write
\[
 c(k)=\int_{\mathbb T}e^{2\pi ikx}\,d\mu(x),\qquad
 E_\mu(N)=\sum_{1\le |k|\le N}|c(k)|^2.
\]
Invariance gives the exact identities $c(2k)=c(3k)=c(k)$. If $\mu$ is not
Haar, some $k_0\ne0$ has $c(k_0)\ne0$: otherwise integration agrees with
Haar on trigonometric polynomials and hence, by uniform approximation,
on every continuous function. Set $K=|k_0|$. Every positive frequency $2^a3^jK$ has the same
absolute coefficient, so
\[
 E_\mu(N)\ge |c(k_0)|^2\,
       \#\{(a,j)\in\Z_{\ge0}^2:2^a3^j\le N/K\}.
\]
The frequencies are distinct by unique prime factorization. If
$L=\log(N/K)$, the lattice count is
\begin{equation}\label{eq183-count}
 \#\{a\log2+j\log3\le L\}
       =\frac{L^2}{2\log2\log3}+O(L+1).
\end{equation}
To justify the error without equidistribution assumptions, tile the
nonnegative quadrant by unit squares anchored at lattice points. The
triangle of weighted radius $L$ is contained in the union of squares
whose anchors lie in that triangle, and that union is contained in the
triangle of weighted radius $L+\log2+\log3$. Their areas differ by
$O(L+1)$, proving \eqref{eq183-count}.

It follows that every non-Haar invariant measure satisfies
\begin{equation}\label{eq183-lower}
 \liminf_{N\to\infty}\frac{E_\mu(N)}{(\log N)^2}
                 \ge\frac{|c(k_0)|^2}{2\log2\log3}>0.
\end{equation}
Thus the following precise estimate would settle the remaining case:
\begin{equation}\label{eq183-missing}
 E_\mu(N)=o((\log N)^2)
 \quad\text{for every nonatomic jointly ergodic invariant }\mu.
\end{equation}
This is a conditional criterion, not an established decay estimate.
Indeed it is equivalent to the desired nonatomic conclusion by the
argument just given and the fact that Haar has $E_\mu(N)=0$. The replication step
itself uses no ergodicity and extends to any multiplicatively independent
integer pair, with the corresponding logarithmic constant.

\paragraph{Attempt: what nonatomicity actually supplies.}
For any probability measure, the elementary Wiener calculation gives
\begin{equation}\label{eq183-wiener}
 \lim_{N\to\infty}\frac{1}{2N+1}
             \sum_{|k|\le N}|c(k)|^2
                        =\sum_x\mu(\{x\})^2.
\end{equation}
Here is the limiting justification. Expand the finite Fourier sum as the
integral over $\mu\times\mu$ of
$D_N(x-y)/(2N+1)$, where
$D_N(t)=\sum_{k=-N}^{N}e^{2\pi ikt}$. Its absolute value is at most one;
it tends to zero for $t\ne0$ in $\mathbb T$ and equals one for $t=0$.
Dominated convergence gives the mass of the diagonal, which is the right
side of \eqref{eq183-wiener}. Hence a nonatomic measure has only
$E_\mu(N)=o(N)$ by this argument. This is compatible with
\eqref{eq183-lower}; dividing by $N$ has discarded the required scale.

To test whether positivity or singularity alone can improve it, take the
usual middle-thirds Cantor measure $\nu$: its ternary digits are independent
and equally likely to be zero or two. It is nonatomic and invariant under
$T_3$. Each of its $2^m$ level-$m$ intervals has length $3^{-m}$ and mass
$2^{-m}$. Let $J=\lfloor3^m/4\rfloor$, for $m\ge2$, and use the
nonnegative Fej\'er kernel
\[
 F_J(t)=\frac1{J+1}\left(\frac{\sin\pi(J+1)t}{\sin\pi t}\right)^2.
\]
For $|t|\le3^{-m}$ the sine inequalities on a fixed subinterval of
$(-\pi/2,\pi/2)$ give $F_J(t)\ge c(J+1)$, with an absolute $c>0$.
Pairs in the same level-$m$ interval have total $\nu\times\nu$ mass
$2^{-m}$. Since $F_J\ge0$ elsewhere,
\[
 \sum_{|k|\le J}\left(1-\frac{|k|}{J+1}\right)|\widehat\nu(k)|^2
   =\iint F_J(x-y)\,d\nu(x)d\nu(y)
   \ge c(J+1)2^{-m}.
\]
The weighted sum is at most $1+E_\nu(J)$. Therefore along these $J$,
\begin{equation}\label{eq183-cantor}
       E_\nu(J)\ge c'J^{1-\log2/\log3}-1.
\end{equation}
This is much larger than $(\log J)^2$, despite nonatomicity, positivity,
and one expanding-map invariance. It is not a counterexample to the
conjecture: $\nu$ is not doubling-invariant. For example $2/9$ is in its
support whereas its double $4/9$ lies in the omitted middle third;
continuity and the definition of support rule out $T_2$-invariance.

\paragraph{Attempt: try to amplify one coefficient.}
For characters $X=e^{2\pi iux}$ and $Y=e^{2\pi ivx}$, Cauchy--Schwarz
applied after subtracting their expectations gives
\begin{equation}\label{eq183-covariance}
 |c(u-v)-c(u)\overline{c(v)}|
          \le\sqrt{(1-|c(u)|^2)(1-|c(v)|^2)}.
\end{equation}
If $u,v$ belong to the same multiplicative orbit of $k_0$, then
$c(u)=c(v)$, and
$\operatorname{Re}c(u-v)\ge2|c(k_0)|^2-1$.
This can produce nontrivial additional coefficients when
$|c(k_0)|>1/\sqrt2$. But a nonzero coefficient need not be that large,
and no amplification to this threshold has been proved. Moreover, taking
all pairwise differences of the $O((\log N)^2)$ replicated frequencies
below $N$ produces at most $O((\log N)^4)$ frequencies. Even persistent
coefficients on that many frequencies do not contradict the $o(N)$ bound.
An iterative argument would need both sustained amplitudes and a
quantitative growth theorem for new frequencies, neither of which follows
from \eqref{eq183-covariance} alone.

\paragraph{Stress test.}
The lower bound has a measure-dependent constant, which is enough to
contradict a little-oh upper bound for that same measure. No uniform lower
bound over all non-Haar measures is asserted. The lattice count concerns
a semigroup of nonnegative powers, not all integer frequencies, and its
zero density is precisely why the Wiener estimate does not settle the
problem. The Fej\'er calculation is exact; no finite simulation has been
promoted to a classification theorem.

The same conditional calculations apply to any coprime $a,b\ge2$, replacing
$2,3$ throughout. This does not prove that solving the single $(2,3)$
problem logically implies the whole parameterized family mentioned in the
catalogue quotation. Both that scope distinction and the original joint
rather than individual ergodicity are retained.

\paragraph{Outcome.}
\textbf{Outcome: Conditional reduction.}

The attempt proves the atomic classification, the logarithmic-square
Fourier obstruction, and two tests explaining why basic Fourier
positivity and nonatomicity estimates do not close the gap. The criterion
\eqref{eq183-missing} makes the missing quantitative input explicit; it
is not a newly proved decay theorem or a claim of priority for Fourier
reformulations.

What remains is an argument using both expanding maps to force that decay,
or another rigidity mechanism in the nonatomic zero-entropy case. The
Cantor construction violates only a weaker one-map proposal, not the
selected conjecture. Neither the full $(2,3)$ classification nor the
broader coprime family is resolved here.
% END RESEARCH ADDITION ATTEMPT-183
\subsection{Sources}
\begin{itemize}\raggedright
\item[\textnormal{[S1]}]\hypertarget{source:183:1}{} P. Burton, J. Panangaden, 'Formulations of Furstenberg's x2 x3 conjecture in complex analysis and operator algebras', arXiv:2410.22701 (2024).
\url{https://ar5iv.labs.arxiv.org/html/2410.22701}.
{\small\textit{Catalogue formulation source.}}
% BEGIN RESEARCH ADDITION SOURCES-183
\item[\textnormal{[E1]}]\hypertarget{extra:183:1}{} Daniel J. Rudolph,
\emph{$\times2$ and $\times3$ invariant measures and entropy},
Ergodic Theory and Dynamical Systems 10(2) (1990), 395--406,
abstract and Sections 1 and 3.
\url{https://doi.org/10.1017/S0143385700005629}.
{\small\textit{Consulted 27 September 2026. The positive-entropy theorem
requires joint ergodicity, not individual ergodicity of both generators.}}
% END RESEARCH ADDITION SOURCES-183
\end{itemize}

\end{document}
